% Generated draft: requires mathematical and source review.
\documentclass[11pt]{article}
\usepackage[margin=25mm]{geometry}
\usepackage{fontspec}
\setmainfont{Latin Modern Roman}
\usepackage{amsmath,amssymb,mathtools,unicode-math}
\setmathfont{Latin Modern Math}
\usepackage{enumitem,needspace,xurl,hyperref}
\setlength{\parindent}{0pt}
\setlength{\parskip}{5pt}
\setlength{\emergencystretch}{3em}
\newcommand{\N}{\mathbb{N}}
\newcommand{\Z}{\mathbb{Z}}
\newcommand{\Q}{\mathbb{Q}}
\newcommand{\R}{\mathbb{R}}
\newcommand{\C}{\mathbb{C}}
\newcommand{\F}{\mathbb{F}}
\newcommand{\A}{\mathbb{A}}
\newcommand{\PP}{\mathbb{P}}
\newcommand{\E}{\mathbb{E}}
\newcommand{\eps}{\varepsilon}
\newcommand{\dd}{\,\mathrm{d}}
\begin{document}
\section*{UnsolvedMath ID: 3100070}
\textbf{Source problem code:} catalogue entry
\subsection{Definitions and mathematical statement}
Let $\varphi(n)$ be Euler's totient function,
\[
 \varphi(n)=\#\{a:1\leq a\leq n,\ \gcd(a,n)=1\},
\]
so $\varphi(1)=1$ and for $n\geq2$ this is the number of positive integers less than $n$ coprime to $n$. Define
\[
 S=\sum_{n=1}^{\infty}\frac{\varphi(n)}{2^n}.
\]
The series converges absolutely since $0\leq\varphi(n)\leq n$ and $\sum n2^{-n}<\infty$. The convention at $n=1$ matters for the numerical value and is fixed by the standard totient function.

\textbf{Questions attributed to Erd\H{o}s.} Determine $S$, with a closed-form answer desired. Is $S$ irrational, that is, $S\notin\Q$? The closed-form request and irrationality question are separate: a numerical approximation or a different infinite-series representation does not by itself provide the requested closed form. The source specifies no class of permitted constants or operations for such a form.

The registry status is open and its classification is OPEN-TRIAGE. The supplied research summary instead claims irrationality known, says no simple elementary closed form is available, and recommends partial-progress wording. It tentatively names Erd\H{o}s--Borwein--Chamayou with a question mark, later names Borwein--Chamayou, and refers to generalized continued-fraction irrational sums and an unspecified average-order expression. It supplies no precise theorem, article, or actual expression verifying those statements. Therefore the irrationality claim remains explicitly attributed and unverified; it is not treated as a settled conclusion of this draft.
\subsection{Short English statement}
Find a closed form for $\sum_{n\geq1}\varphi(n)/2^n$ and determine whether it is irrational. The standard convention is $\varphi(1)=1$. The summary claims irrationality known but gives tentative attribution and no verifiable theorem reference, conflicting with the open registry status.
\subsection{Sources}
\begin{itemize}
\item Local UnsolvedMath registry, permanent ID 3100070, source code catalogue entry, title beginning ``What is the sum of phi(n)/2 to the n''; proposer: Erd\H{o}s. Supplied status open, classification OPEN-TRIAGE, updated 2026-08-06; the research summary separately recommends partial-progress wording.
\item Cooper, \emph{Combinatorial Problems I Like}, October 2020 snapshot, section ``Classical Number Theory,'' source-order bullet 70, \url{https://people.math.sc.edu/cooper/combprob.html}. The record reports HTML extraction and access on 2026-07-29; later resolution was not checked item by item. Rights: NEEDS\_REVIEW; public access verified but redistribution permission not established.
\item The supplied summary mentions Erd\H{o}s, Borwein, and Chamayou with expressly uncertain attribution and no paper title or theorem; it supplies no actual average-order formula. No bibliographic details are reconstructed.
\end{itemize}
\end{document}
